\documentclass[11pt]{article}
\usepackage[letterpaper,margin=0.85in]{geometry}\usepackage[T1]{fontenc}\usepackage[utf8]{inputenc}\usepackage{amsmath,amssymb,enumitem,booktabs,array,xcolor,fancyhdr,tabularx}
\setlist{itemsep=2pt,topsep=3pt}\pagestyle{fancy}\fancyhf{}\renewcommand{\headrulewidth}{0.3pt}
\fancyhead[L]{\small Effective Nuclear Charge and Shielding}\fancyhead[R]{\small Make-up key}\fancyfoot[C]{\small\thepage}
\definecolor{ans}{RGB}{28,125,74}\newcommand{\A}[1]{\textcolor{ans}{\textbf{#1}}}
\begin{document}

\section*{Effective Nuclear Charge and Shielding}
\textbf{You missed this lesson. Copy this into your notes. That is all you need to do.}\par\medskip
\subsection*{Question of the day}
How hard does the nucleus really pull on an outer electron?\par
\subsection*{Targets}\begin{itemize}\item I can estimate Z\_eff as Z minus core electrons and use it to compare how hard atoms hold electrons.\end{itemize}
\subsection*{Concepts}\begin{itemize}
\item \textbf{Core electrons} (core = Z -- outer electrons): Every electron that is not in the outer ring.
\item \textbf{Shielding} (a core electron cancels one proton): Core electrons cancel part of the nucleus's pull.
\item \textbf{Z\_eff} (Z\_eff = Z -- core): The pull an outer electron really feels. Estimate: Z minus the core electrons.
\item \textbf{Distance} (pull = Z\_eff $\div$ n\textsuperscript{2}): The farther the ring, the weaker the pull (slot 7).
\end{itemize}
\subsection*{What the class predicted, and what is true}\begin{itemize}
\item Lithium and sodium each have one outer electron. Which atom holds its outer electron harder? \A{Lithium.} Sodium's outer electron is in ring 3, lithium's in ring 2. The farther electron is pulled less.
\end{itemize}
\section*{Huddle notes (answers filled in)}
\subsection*{Definitions}\begin{itemize}
\item \textbf{Core electrons:} The electrons \A{inside} the outer ring.
\item \textbf{Shielding:} Core electrons \A{cancel} part of the nucleus's pull.
\item \textbf{Z\_eff:} Z\_eff = Z minus the \A{core} electrons.
\end{itemize}
\subsection*{Worked example}
Sodium (Z = 11): how hard is its outer electron held?\\
\begin{itemize}
\item Rings: 2, 8, 1. Total \A{11}.
\item Core = 2 + 8 = \A{10}.
\item Z\_eff = 11 -- 10 = \A{1}.
\item Pull = 1 $\div$ 3\textsuperscript{2} = \A{1/9}.
\end{itemize}
\subsection*{Label the sodium atom}
\begin{enumerate}
\item The nucleus has charge +\A{11}.
\item The \A{core} rings hold \A{10} electrons.
\item The \A{outer} ring holds \A{1} electron.
\end{enumerate}
\section*{One worked sample of each skill}
\subsection*{Count the core electrons}
\noindent\textbf{Level 1.} How many core electrons does this sodium atom have? Core electrons are the ones not in the outer ring.\\ \A{Answer: 10} \quad Rings: 2, 8, 1. The outer ring is not core, so 2 + 8 = 10.\par\medskip
\noindent\textbf{Level 2.} A neutral magnesium atom has Z = 12 and its electrons sit in rings of 2, 8, 2. How many are core electrons?\\ \A{Answer: 10} \quad The last ring (2) is the outer ring. Core: 2 + 8 = 10.\par\medskip
\noindent\textbf{Level 3.} A neutral atom has Z = 19 and 1 electron in its outer ring. How many core electrons does it have?\\ \A{Answer: 18} \quad 19 -- 1 = 18 core electrons.\par\medskip
\noindent\textbf{Level 4.} A neutral atom has Z = 18 and 8 electrons in its outer ring. How many core electrons does it have?\\ \A{Answer: 10} \quad 18 -- 8 = 10 core electrons.\par\medskip
\subsection*{Estimate Z\_eff}
\noindent\textbf{Level 1.} The picture shows sodium (Z = 11). Z\_eff is Z minus the core electrons. What is Z\_eff for sodium?\\ \A{Answer: 1} \quad 11 -- 10 = 1.\par\medskip
\noindent\textbf{Level 2.} Estimate Z\_eff for sodium (Z = 11, rings of 2, 8, 1).\\ \A{Answer: 1} \quad Core = 10. Z\_eff = 11 -- 10 = 1.\par\medskip
\noindent\textbf{Level 3.} An atom in period 3 has Z\_eff = 1. Which atom is it?\\ (A) Phosphorus \quad \A{(B) Sodium} \quad (C) Chlorine \quad (D) Aluminum\\ \A{Answer: B} \quad Z\_eff = Z -- 10, so Z = 1 + 10 = 11: sodium.\par\medskip
\noindent\textbf{Level 4.} An atom in period 3 has Z\_eff = 8. Which atom is it?\\ \A{(A) Argon} \quad (B) Silicon \quad (C) Magnesium \quad (D) Chlorine\\ \A{Answer: A} \quad Z\_eff = Z -- 10, so Z = 8 + 10 = 18: argon.\par\medskip
\subsection*{Which holds its outer electron harder}
\noindent\textbf{Level 1.} Each atom has one outer electron. Which atom holds its outer electron harder, lithium (left) or potassium (right)?\\ \A{(A) Lithium} \quad (B) Potassium \quad (C) They hold it equally\\ \A{Answer: A} \quad Both have Z\_eff = 1. Lithium's outer electron is in ring 2 and potassium's is in ring 4, so lithium pulls harder.\par\medskip
\noindent\textbf{Level 2.} Which atom holds its outer electron harder, calcium (Z = 20) or magnesium (Z = 12)? Find Z\_eff for each.\\ (A) Calcium \quad \A{(B) Magnesium} \quad (C) They hold it equally\\ \A{Answer: B} \quad Both have Z\_eff = 2. Magnesium's outer electron is in ring 3 and calcium's is in ring 4, so magnesium pulls harder.\par\medskip
\noindent\textbf{Level 3.} Which atom holds its outer electron harder, silicon (Z = 14) or calcium (Z = 20)? Use Z\_eff and the outer ring number.\\ (A) Calcium \quad \A{(B) Silicon} \quad (C) They hold it equally\\ \A{Answer: B} \quad Silicon: 4 $\div$ 3\textsuperscript{2} = 4/9. Calcium: 2 $\div$ 4\textsuperscript{2} = 2/16. The larger number pulls harder.\par\medskip
\noindent\textbf{Level 4.} Which atom holds its outer electron harder, argon (Z = 18) or silicon (Z = 14)? Use Z\_eff and the outer ring number.\\ (A) They hold it equally \quad (B) Silicon \quad \A{(C) Argon}\\ \A{Answer: C} \quad Both outer electrons are in ring 3. Argon has Z\_eff = 8 and silicon has 4, so argon pulls harder.\par\medskip
\end{document}